\VerdierPhysicalPageAnchor{0049}%
Pour tout élément $\sigma\in G$, le diagramme :

\VerdierSemanticAnchor{ch1-1-4-5-eq4}
\begin{center}
  \resizebox{0.78\textwidth}{!}{%
    $\begin{tikzcd}[ampersand replacement=\&,column sep=6em,row sep=4em]
      F_1T(\sigma)
        \arrow[r,"{p_1(\sigma)}"]
        \arrow[d,"{u*T(\sigma)}"']
        \& T'(\sigma)F_1
          \arrow[d,"{T'(\sigma)*u}"] \\
      F_2T(\sigma)
        \arrow[r,"{p_2(\sigma)}"]
        \& T'(\sigma)F_2
    \end{tikzcd}$%
  }
  \hfill\textup{(1.4.5.4)}
\end{center}
est $\beta(\sigma)$-commutatif.

\VerdierSemanticAnchor{ch1-1-4-6}
\noindent\textbf{Définition 1.4.6.}\quad
Soient $(\mathcal C,G,T,c)$ une $G$-catégorie,
$(\mathcal C',G,T',c')$ une $G$-catégorie additive,
$\alpha_1$ et $\alpha_2$ dans $Z^2(G,\mathbb Z/2\mathbb Z)$ deux
cocycles non homogènes normalisés,
$\beta\colon G\to\mathbb Z/2\mathbb Z$ une application telle que :

\VerdierSemanticAnchor{ch1-1-4-6-eq1}
\begin{equation}
  \alpha_2(\sigma_1,\sigma_2)\alpha_1^{-1}(\sigma_1,\sigma_2)
  =\beta(\sigma_2)\beta(\sigma_1\sigma_2)^{-1}\beta(\sigma_1).
  \tag{1.4.6.1}
\end{equation}

\VerdierSemanticAnchor{ch1-1-4-6-beta-morphism}
Soient enfin
$(F_1,p_1),(F_2,p_2)\colon(\mathcal C,G,T,c)
\to(\mathcal C',G,T',c')$ deux $G$-foncteurs tordus respectivement
par les cocycles $\alpha_1$ et $\alpha_2$. On appelle
\emph{$\beta$-morphisme du foncteur $(F_1,p_1)$ dans le foncteur
$(F_2,p_2)$} un morphisme $u\colon F_1\to F_2$ de foncteurs tel que,
pour tout $\sigma\in G$, le diagramme (1.4.5.4) soit
$\beta(\sigma)$-commutatif.

\VerdierSemanticAnchor{ch1-1-4-7}
\noindent\textbf{1.4.7.}\quad
On a montré (1.4.5) que le choix de deux sections normalisées permet
d'associer à un morphisme de $H$-foncteurs induisant l'opération
canonique du groupe des signes un $\beta$-morphisme entre les
$G$-foncteurs tordus correspondants. On vérifie immédiatement que
cette correspondance est biunivoque. De plus, lorsqu'on a trois
$H$-foncteurs $(F_1,m_1)$, $(F_2,m_2)$, $(F_3,m_3)$ induisant
l'opération canonique du groupe des signes, trois sections normalisées
$s_1$, $s_2$, $s_3$ et deux morphismes de $H$-foncteurs $u$, $v$ et
leur composé $vu$ :

\VerdierSemanticAnchor{ch1-1-4-7-eq1}
\begin{equation}
  \begin{aligned}
    u&\colon(F_1,m_1)\longrightarrow(F_2,m_2),\\
    v&\colon(F_2,m_2)\longrightarrow(F_3,m_3),\\
    vu&\colon(F_1,m_1)\longrightarrow(F_3,m_3),
  \end{aligned}
  \tag{1.4.7.1}
\end{equation}
les morphismes :

\VerdierSemanticAnchor{ch1-1-4-7-eq2}
\begin{equation}
  \begin{aligned}
    u&\colon(F_1,p_1)\longrightarrow(F_2,p_2),\\
    v&\colon(F_2,p_2)\longrightarrow(F_3,p_3),\\
    vu&\colon(F_1,p_1)\longrightarrow(F_3,p_3),
  \end{aligned}
  \tag{1.4.7.2}
\end{equation}

\VerdierSemanticAnchor{ch1-1-4-7-open-seam}
